


Date: 2026-11-13. Branch \texttt{finish} (from tag \texttt{positive-push-v1}). Software only; \texttt{data/hardware} is not read. Result keys start with \texttt{f.}; no tagged result and no existing key is changed. Nothing is pushed or uploaded.
What was seen before this text: every result of the earlier rounds (V1 to V4, UA-SCT, X1 to X8, P1 to P4d, T-series). What was NOT seen: any number of Blocks B, C, F or G. Anything not written here is EXPLORATORY and says so. A change after this commit is a dated, labelled deviation in the addendum at the end and in \texttt{docs/07\_log/DECISIONS.md}; sentences above are never edited.
Labels: MEASURED\_SOFTWARE, SIMULATED, ANALYTICAL, NOT\_RUN. Anything not completed is recorded NOT\_RUN with its reason; this is allowed and is not a deviation. Wording rule of the repository applies (the words forbidden by CLAUDE.md, rule 7, are not used).

\subsection*{0. Common rules}

\begin{itemize}
\item \textbf{Revised after the negative.} Every B2, F and G item that revisits an earlier negative result is a revised design tested after seeing the negative. It is reported next to the original result, never in place of it, and tables say "revised after the negative". Confirmation counts only on scenarios, seeds and records not used to design it. F results are labelled EXPLORATORY-REVISED everywhere.
\item \textbf{Seeds.} 42 for single-seeded work; 0 to 4 for repeated runs of models with randomness. Simulation seeds (hour index h = 0 to 29, 30 seeds per cell, always paired across policies): existing design scenarios of the UA-SCT study use seed 3000 + h (the stored runs); B2 held-out 5000 + h; F3 9000 + h; G DESIGN grid 6000 + h; G HELD-OUT grid 8000 + h. Synthetic pulse generator seed 42 (pulse k uses seed 42 + k).
\item \textbf{Intervals.} 95\% percentile bootstrap, 1,000 resamples. Simulation cells: seeds resampled inside each cell; across-cell statements count cells. NHTSA records: whole make-and-model groups resampled together; paired across compared methods. VZCrash: test events resampled; paired.
\item \textbf{A cell "gains" if} the point difference in the stated metric is at least the stated margin AND the 95\% interval excludes zero. A cell "loses by more than 2 points" if the point difference is below -0.02. (Applies to H-B2, H-F3a, H-F3b, H-G3a, H-G3b, H-G3c.)
\item \textbf{Holm correction} over the 18 hypotheses of this document: H-B2, H-C3, H-F1, H-F1b, H-F1c, H-F1d, H-F1e, H-F2a, H-F2b, H-F2c, H-F2d, H-F2e, H-F2f, H-F3a, H-F3b, H-G3a, H-G3b, H-G3c. Each has one p-value: the largest one-sided bootstrap p-value (share of resamples on the wrong side of the criterion; floor 0.001) among the conditions that must all hold; for "in at least a share q of cells", the p-value is the one of the cell at the q-quantile of the per-cell p-values (cells ordered from smallest p). SUPPORTED = criterion met and Holm-adjusted p below 0.05; MIXED = criterion met, adjusted p not below 0.05; NOT SUPPORTED otherwise. F hypotheses keep that verdict word and carry the extra label EXPLORATORY-REVISED.
\item \textbf{Data origin and splits.} NHTSA pool: the 913 records and the make-and-model group folds of \texttt{docs/05\_data\_code/SEALED\_SPLIT.md} (development 543, validation 184, sealed 186). VZCrash: supplied splits; primary test set de-duplicated (24,022 events), strict set 22,830 events; thresholds from validation only; test touched once per model and seed (revised methods are scored on the already-scored sealed fold and test split once each, labelled "revised after the negative"; the main estimate is grouped nested cross-validation, rule 11).
\item \textbf{Not a real-world claim.} Simulated results are statements about the simulator under numbered assumptions (section G0). Synthetic pulses are a stress test; they are never called unbiased and never a substitute for real data.
\item \textbf{Sources.} No paper, datasheet or dataset is cited unless opened; URL and date are recorded in \texttt{docs/06\_references/REFERENCES.md}; a source that cannot be opened is said so and none of its numbers is used.
\end{itemize}

\subsection*{B. Problems found in simulated and protocol work}

\textbf{B1 (ANALYTICAL; a check of the model within bounds, not of the implementation).} Property C3 of \texttt{protocol\_model\_ext} (a CANCEL that was delivered is never lost because it overtook its alert) fails for the design variant in 351,754 reachable states (bounds EXT, the T12 trace: alert and CANCEL sent, CANCEL delivered ahead of the alert and refused as unknown, alert delivered later, event active although cancelled). Proposed minimal change: the base remembers a valid CANCEL for a counter it has not seen yet (exact counter, bounded memory) and applies it when that alert arrives (\texttt{keep\_pending\_cancel} of T12). Re-run with the same bounds EXT and all original properties (S1 to S4, C1 to C4, F1). Reported: states explored, depth, properties checked, whether the counterexample disappears. A watermark form (a single cancelled-up-to counter) is NOT proposed: it would cancel a delayed earlier event; an exploratory two-event bound shows this if it completes. The reference code (\texttt{sct\_protocol.SctBase}) gets an option \texttt{pending\_cancel} (default off: the 48 + 9 attack cases of the tagged code are unchanged); with it on, every attack case is re-run and X06 is expected to change from "not cancelled" to "cancelled"; every other case is expected unchanged. Not a hypothesis.

\textbf{B2 (SIMULATED). Which parts of the fast path carry the gain.} Earlier findings: ablations at +-16 g found that flash and hold carry the gain; the equal-escalation test (N28) found the interval adds nothing over a lowered table. Simple consequence tested here: \textbf{P1+} = the severity-table policy (P1: tries 2, 4, 8, 12 by class, service by class, immediate retry) plus the two parts: (a) flash ahead of the alert for Severe and Critical (14-byte flash, then the 97-byte alert) and (b) hold of 0 s for Critical and \texttt{s11.tuned.h\_s} s for Severe (Low and Moderate keep the 10 s and 8 s windows). No SF raise, no reserved duty share, no pre-emption, no spaced retries. Comparators: P0 (fixed 3 retries), P1 (current thresholds), \textbf{P1e} (P1 with every class boundary divided by the factor f of \texttt{t1.r16.f}, chosen on real pool rows to equal the escalation of P3), \textbf{P1+e} (P1+ with the same lowered classes), P3 (UA-SCT) for reference. Class pools: the +-16 g NHTSA cross-validated pairs of the UA-SCT study. Scenario: one cell has n cars, packet loss p (up and down), SF, a burst of b Critical events at 600 s plus background events, duty 1\%, hb 600 s, 16 byte heartbeats (all other values as the UA-SCT study).
\begin{itemize}
\item DESIGN: the existing UA-SCT grid at +-16 g (9 burst cells, 7 steady cells), seeds 3000 + h, h = 0 to 29. P0 and P1 are the stored runs; P1+ and P1+e are new runs on the same seeds. Used only to look at the size and direction of the effect; no parameter is tuned (P1+ uses the stored tuned values).
\item \textbf{HELD-OUT (confirmation):} cars \{75, 150\} x packet loss \{10\%, 30\%, 50\%\} x SF 8 x burst size \{20, 45\} = 12 cells, +-16 g, 30 fresh seeds per cell (5000 + h).
\item \textbf{H-B2.} P1+ beats P0 on critical reach within 10 s by at least 5 points with the interval excluding zero in at least 60\% of the held-out cells (at least 8 of 12), AND loses by more than 2 points in none. \textbf{Counts against:} fewer than 60\% of cells, or any cell lost by more than 2 points. p-value rule 5 (cell at the 60\% quantile; and the largest per-cell p of the losing clause).
\item Also reported (descriptive, with intervals): P1+ minus P1, P1+e minus P1e (flash and hold at equal escalation: identical class assignment, so escalation is equal by construction), P1+e minus P3, critical reach within 30 s, critical p95 latency, duty-cycle use (node airtime share and fast-path share), the share of non-severe events escalated (\texttt{escalation\_used}), duty violations (must be 0). +-32 and +-64 g cells are EXPLORATORY.
\end{itemize}

\textbf{B3 (ANALYTICAL and SIMULATED). Simulator sanity checks.} (i) Unit tests of \texttt{lora.airtime\_ms} against Semtech's published time-on-air formula for SF7 to SF9, 125 kHz, coding rate 4/5, 8-symbol preamble, explicit header, CRC on, payloads 14, 16, 26 and 97 bytes; the document opened is cited with URL and date; if it cannot be opened the test is recorded as against the formula as transcribed and says so. (ii) Periodic traffic with no capture (\texttt{capture="none"}, no alerts, no extra loss, SF7, 16-byte heartbeats, 1,000 s period jittered by the simulator's own rule, n cars chosen to span offered load G = n x airtime / period from 0.01 to 1): the simulated collision probability of an unconfirmed heartbeat against 1 - exp(-2G); the error (absolute and relative) is reported; mismatches are limitations and the simulator is not tuned. (iii) The duty-cycle limit (independent scan of logged transmissions) is never exceeded in any run of any cell reported by this document; the count of runs is reported and must be 0.

\textbf{B4.} \texttt{docs/03\_results/NEGATIVES\_REGISTER.md}: one row per negative result (N20, N28, H-P1 to H-P4c, H-X1 to H-X8, T4 single-car loss case, P3b), original verdicts unchanged, inferred causes marked "inferred". Not a hypothesis.

\subsection*{C. Synthetic stress-test pulses (SIMULATED)}

\textbf{C1 generator.} Per crash mode (frontal barrier, pole, side, other) a pulse-family generator on the horizontal plane at 10 kHz: families half-sine, triangle (peak at fraction a of the duration) and multi-peak (2 or 3 half-sine bumps); the family is chosen with probability 1/3 each (assumption A-C1). Parameters per mode: log peak resultant acceleration, log duration, logit of the asymmetry a, drawn from a multivariate normal fitted to the pooled real pulses of that mode (mean and covariance only; all 913 records for the stress-test file and the C-block results; the development fold only, refitted inside each training fold, for any F-block training use, so that no validation or sealed information enters training). Direction angle drawn from the empirical angle distribution of the mode. Additive band-limited noise with relative level fitted to the real pulses. Quiet pre-trigger 0.1 s. Every numbered assumption (value and source) is a line in \texttt{docs/05\_data\_code/SYNTHETIC\_ASSUMPTIONS.md}. Published: code, seeds and a fixed file of 5,000 pulses (parameters and 2 kHz waveforms; modes in the proportion of the pool).
\textbf{C2 fidelity.} Per mode, real against synthetic: two-sample Kolmogorov-Smirnov on peak resultant g, duration (width at 10\% of the peak), delta-V (the N13 reference), asymmetry (time of peak over duration) and a shape descriptor (area over peak x duration); plots. \textbf{A mode passes if the KS statistic is at most 0.20 for every one of the five descriptors.} A failing mode is excluded from C3, F and G and this is said. The fit is in-sample for the C-block (stated).
\textbf{C3 use.} The N13 sensor-range grid (B2 of V4: eight real parts, ranges, rates, plain delta-V within +-10\% of the reference, rule: lower end of the 95\% interval of the share at least 0.80), run on the synthetic pulses of the modes that passed (independent pulses: 1,000 per mode, bootstrap over pulses).
\begin{itemize}
\item \textbf{H-C3.} For each overlapping mode that passed fidelity (frontal barrier, pole), the smallest range meeting the 0.80 rule on the synthetic pulses is within one range step of the real-data result (\texttt{p.b2.<mode>.smallest.config}). A range step is one rung of the ladder of ranges of the eight parts (sorted union of \texttt{sel\_parts.study\_ranges}). \textbf{Counts against:} it differs by more than one range step for any overlapping mode that passed. If no mode passes: NOT\_RUN. p-value: not defined for an exact-match criterion; H-C3 enters Holm with p = 0.001 if it holds in all passing modes and 1.0 otherwise (stated; a criterion about agreement has no sampling null).
\item Extension (EXPLORATORY, assumption-dependent): the log-peak mean of each mode shifted by +0.2, +0.4 and +0.6 (A-C2) to represent more severe crashes; required range reported per shift; sensitivity of the required range to each assumption varied one at a time (family mix 1/3 each against 0.5/0.25/0.25 and 0.25/0.25/0.5; duration scale -20\% and +20\%; correlation of log peak and log duration +-0.2; noise level x0.5 and x2; pre-trigger length). Every figure and table from synthetic pulses is captioned "synthetic stress test, assumptions in SYNTHETIC\_ASSUMPTIONS.md".
\end{itemize}

\subsection*{F. Reruns of the failed algorithms with real plus synthetic training data (all EXPLORATORY-REVISED)}

\textbf{F rules.} (1) Synthetic data train only. Validation, conformal calibration, threshold and ratio selection and every reported test metric use real data only; no synthetic record is in any score called a real-data result. (2) Gate: a synthetic mode or family is used only if it passed C2; the used set is recorded. (3) Mixing ratio r = synthetic records per real training record, chosen on the REAL validation fold from \{0, 0.5, 1, 2, 5\}; r = 0 is the control and uses identical code (same augmentation, same pipeline), so the effect of synthetic data is separated from the effect of new code. (4) Per task: the original result, the real-only rerun (r = 0), the best r, an interval; a statement whether any gain exists only at r > 0. (5) TSTR diagnostic for each model family: train on synthetic only, test on real (a fidelity measure; never a result, and a real-to-synthetic score is never reported). (6) Original criteria and thresholds of each hypothesis are reused unchanged. (7) The sealed NHTSA fold and the supplied VZCrash test split were scored by the original methods; each revised method is scored on them once, labelled "revised after the negative"; the main estimate is grouped nested cross-validation over all real records (groups together, stratified by mode on the NHTSA side; 5 outer folds, the inner validation fold chosen group-wise inside the training part), with Holm over the F hypotheses (rule 5). (8) Nothing about the sealed fold or test split is used to choose r or any setting.
\textbf{F0 VZCrash-style surrogate generator.} Accelerometer windows (280 samples x 3 axes at 100 Hz, trigger at sample 100), crash and non-crash, fitted to SUMMARY statistics of the training split only (class balance, peak and duration distributions, spectral band shares, onset behaviour (pre-trigger level and rise time)), never to individual records; only fitted parameters and code are published. Validation: marginals and spectra against the supplied windows by two-sample tests and plots; nearest-neighbour check of every synthetic window against all real training windows (L-infinity on the 32-bin downsampled |a| and on the full window) with the distance distribution reported against the real-to-real near-duplicate distances; no synthetic window may be a copy (minimum distance above the 1st percentile of real-to-real nearest-neighbour distances, else the offending windows are dropped and counted). README states: VZCrash is CC BY-NC 4.0 (Hugging Face revision 42e1482d005c0e545379d6cfa4939486c9fe4249); the generator parameters are therefore released under CC BY-NC 4.0 with attribution; only summary parameters are released; it matches only the marginals that were validated; detector scores on synthetic windows are not comparable with real-data scores; the surrogate is not called unbiased.
\textbf{F1 NHTSA side} (913 pool; ranges +-16, +-32, +-64 g, 100 Hz primary): training on development real records (+ augmentation of the original design, 6 copies) plus r x (number of real training records) synthetic pulses passing through the same sensor chain.
\begin{itemize}
\item \textbf{H-F1 (original H-P1).} Best deployable direct model (candidates: augmented direct GBM and augmented 1D CNN, mode none; chosen on validation macro-F1 over model kind and r) against the best of plain, fill + ridge, anchored (original, real-only; the same estimators also trained with the chosen r are reported as the fairer baseline, EXPLORATORY-REVISED): macro-F1 higher by at least 0.03 with the interval above zero AND under-triage lower by at least 0.03 with the interval below zero, in at least two of three ranges. \textbf{Counts against:} in two or more ranges either interval includes zero or the point difference is under 0.03.
\item \textbf{H-F1b (original H-P2).} Interval [plain, anchored estimate x q]; q the conformal quantile of reference / anchored on the REAL validation fold with the V4 alpha rule (target 90\%); anchored trained with r. Coverage lower end at least 0.85 overall and for pole on the sealed fold, AND the relative reduction of the mean upper-to-lower ratio against the existing relative-residual interval around fill + ridge at least 0.20 with the lower end above zero. \textbf{Counts against:} either fails.
\item \textbf{H-F1c (original H-P3a).} Reliability score regressor trained with r; abstain above the 70th percentile of the score on the REAL validation fold. Retained under-triage at most half the unselected rate (ratio at most 0.5 with the upper end below 1) AND retained p90 error at most 20\%.
\item \textbf{H-F1d (original H-X6, splits G and O).} On the 176 original tests (grouped 5-fold, 10 repeats of \texttt{saturation\_splits}), +-16 g, 100 Hz: the gain in p90 absolute error of the anchored estimator over fill + ridge is at least 5 points on split R, and its whole interval above zero on at least one of the splits G, A, O, with under-estimation share below that of the tail-aware estimator. The rerun reports splits G and O (where the original gain failed) and R; the original criterion applied unchanged to the three splits G, O and R.
\item \textbf{H-F1e (original H-X5).} Spearman correlation of fill + ridge delta-V with the reference at least 0.90 and lower end at least 0.85 at +-16 g, 100 Hz, trained with r inside the grouped folds.
\end{itemize}
\textbf{F2 VZCrash side} (5 seeds; primary set de-duplicated; strict set sensitivity): the 15-tree depth-6 forest trained on real triggered training windows plus r x (real count) surrogate windows (features recomputed by the 14-feature code), r chosen on the real validation split (AP).
\begin{itemize}
\item \textbf{H-F2a (original H-P4b).} Recall at the 1\% false-alarm target on the strict set at least 0.02 above standard training (r = 0, real only), interval above zero.
\item \textbf{H-F2b (original H-X3).} Group-aware training (down-weight 1 / group size of real windows; surrogate windows weight 1) raises recall at the 1\% target on the strict set by at least 0.01 over standard real-only training, interval above zero.
\item \textbf{H-F2c (original H-X2).} Trigger cascade: the rule family and selection of V3 use whole 16 s records; surrogate windows are trigger-aligned windows and cannot train or test a trigger. Preregistered as NOT\_RUN unless a record-level surrogate exists (none is planned); reported with that reason, original verdict unchanged.
\item \textbf{H-F2d (original H-X1).} 16-bit forest trained with r: AP loss upper bound at most 0.005 AND flash estimate cut at least 30\% under the primary layout. The flash clause is arithmetic (22.2\% for 16-bit thresholds in the primary layout) and cannot change with training data; the verdict is therefore expected to stay NOT SUPPORTED and the AP clause is reported.
\item \textbf{H-F2e (original H-X8).} Sequential rule (six length-specific forests trained with r) earlier median decision than the early-exit rule by at least 100 ms at equal or better recall, same targets.
\item \textbf{H-F2f (original H-P4c).} Calibration stays on real held-out vehicle-size categories (leave-one-category-out); with the better-trained model, some variant has false-alarm upper end at most alpha + 0.01 at both 1\% and 5\% AND recall at least 0.02 above the trivial rule.
\end{itemize}
\textbf{F3 simulator} (SIMULATED): the improved severity estimates of F1 (pairs from the sealed and nested-CV predictions at +-16 g) feed the UA-SCT simulator on the HELD-OUT grid of B2 (12 cells, seeds 9000 + h, 30 seeds). \textbf{H-F3a (original N28 decision rule):} P3 with the improved estimates beats the lowered-threshold table at equal escalation: pool-level mis-triage interval (P3 minus P1e) entirely below zero AND critical reach within 10 s not worse than -0.02 at the lower end in steady and burst cells. For the held-out grid (burst cells): in at least 60\% of cells P3 gains at least 5 points over P1e or P2e on critical reach within 10 s and no cell loses more than 2 points. \textbf{H-F3b (original H-P3b):} mis-triage of abstain-routing minus lowered table at most -0.02 with the interval below zero, with improved estimates.
\textbf{F4.} Summary table, one row per failed hypothesis: original verdict, real-only rerun, real plus synthetic rerun with best r and interval, TSTR score, verdict, whether the verdict changed. If no verdict changes, this is said plainly.

\subsection*{G. Protocol layers on assumption-driven simulated data (all SIMULATED)}

\textbf{G0 numbered assumptions} (values and sources; also in \texttt{docs/05\_data\_code/SYNTHETIC\_ASSUMPTIONS.md}). Source "assumed" = no document opened; the simulator is \texttt{crashguard.world\_sim} (a copy of \texttt{sct\_sim} with optional world features; with every feature off it reproduces \texttt{sct\_sim} exactly, tested).
\begin{itemize}
\item A-G1 Crash arrivals: Poisson 0.1 per car per hour (stress parameter, as in the UA-SCT study; assumed). Burst: b Critical events within 10 s at 600 s (b in the grids).
\item A-G2 Severity mix: (true class, assigned class) pairs of synthetic pulses of the modes that passed C2, in the mode proportion of the pool; assigned class from the plain estimate at +-16 g, 100 Hz (policies P0 to P2), from the upper end plain x q with q the real validation conformal ratio (P3, P4); the F1-improved estimate where F1 completed. Burst events are drawn from true-Critical entries.
\item A-G3 Vehicle density: n cars share one gateway cell (grid values).
\item A-G4 Geometry and path loss: cars uniform in a disc of radius 3 km; one gateway at the centre (grids) or two gateways at +-1.5 km (G5 only); path loss PL(d) = 25 dB + 10 x 3.0 x log10(d / 1 m) (exponent 3.0 assumed; 25 dB at 1 m assumed); transmit power 14 dBm and 10 dB of vehicle and body loss (assumed); static log-normal shadowing sigma 6 dB per car and gateway (assumed) plus the existing per-transmission fading of 2 dB; reception needs a level above the SF sensitivity (assumed typical values for 125 kHz: SF7 -123, SF8 -126, SF9 -129, SF10 -132 dBm; replaced by a document value if Semtech's datasheet is opened for B3, with its URL and date).
\item A-G5 Capture threshold 6 dB at equal SF (as the UA-SCT study; assumed); other SF orthogonal with the loss factor of the UA-SCT study.
\item A-G6 Independent extra packet loss p on uplink and downlink (grid values).
\item A-G7 SF and coding rate: SF 7 to 10; coding rate 4/5 (G2 varies 4/5 to 4/8); 125 kHz, 8-symbol preamble, explicit header, CRC on.
\item A-G8 Duty-cycle limit 1\% and 0.1\% of any trailing hour for each transmitter including the gateway; \textbf{the India limit is UNVERIFIED} (WPC/DoT not checked).
\item A-G9 ACK timing: the ACK starts 0.1 s after the end of the uplink (assumed; G5 varies it to 0.5 and 1.0 s, a LoRaWAN-like receive window); the car listens for 0.1 s + ACK airtime + margin 0.1 s.
\item A-G10 Channels: 1 (grids); 3 equal channels (G5): each transmission picks one uniformly; collisions only on the same channel; the gateway receives on all channels and is half-duplex on the channel it transmits on (assumed).
\item A-G11 Tag length: 8 bytes (grids); 4 and 16 in G2 and G5 change the lengths of alert, flash, CANCEL and ACK by the difference (assumed layout otherwise as \texttt{docs/01\_protocol/PACKET\_FORMAT.md}).
\item A-G12 Slotted fast path (G3a): the gateway sends a 9-byte beacon every 30 s; a frame is 30 s; each frame starts with a reserved window of 8 slots of length s = max(0.5, ceil(alert airtime / 0.1) x 0.1 + 0.1) s; Severe and Critical fast-path transmissions (flash, alert, CANCEL) start at the beginning of a uniformly chosen slot of the next reserved window; normal-path and heartbeat transmissions avoid the reserved window; a car is synchronised if it heard a beacon in the last 90 s (else it behaves as in the baseline). Beacon airtime is charged to the gateway's duty and reported. Beacon coordination is assumed.
\item A-G13 Listen-before-talk (G3b): before every fast-path transmission the car senses; each ongoing transmission on its channel is detected with probability 0.7 (hidden terminals 0.3; assumed); if busy it waits U(0.05, 0.25) s and senses again, at most 4 times, then transmits. The regulatory status of listen-before-talk at 433 MHz is UNVERIFIED.
\item A-G14 V2V relay (G3c): K = 3 other cars each hear an alert transmission with probability 0.5 (assumed); if no ACK reached them within the listen window of the sender, each hearing car repeats the alert once after U(1, 5) s (assumed); the relayed copy is identical (same counter), is charged to the relaying car's duty, is received by the gateway under the relayer's link, and the gateway re-ACKs idempotently; it is a duplicate when the gateway already holds the alert.
\item A-G15 Receivers: the gateway decodes per gateway independently; with two gateways an uplink is received if any gateway receives it and the ACK goes through the gateway with the best level (G5).
\end{itemize}
\textbf{DESIGN grid (16 cells, fixed here):} n \{75, 150\} x p \{0.1, 0.3\} x SF \{7, 9\} x duty \{1\%, 0.1\%\}; burst b = 20; one gateway; one channel; tag 8 bytes; +-16 g pairs; seeds 6000 + h.
\textbf{HELD-OUT grid (16 cells, fixed here, nothing is chosen on it):} n \{100, 200\} x p \{0.2, 0.4\} x SF \{8, 10\} x duty \{1\%, 0.1\%\}; burst b = 30; same other settings; seeds 8000 + h.
\textbf{G1 policies} on both grids, 30 seeds per cell: P0, P1, P1e, P1+, P2 (SCT), P3 (UA-SCT), P4 (capped; phi = \texttt{s11.tuned.phi}), P3b abstain-routing (A of V4: abstaining events sent to the fast path), and the policies fed by F1 estimates (P1+ and P3 with improved classes) where F1 completed. Metrics: mis-triage, critical reach within 10 s and 30 s, critical p95 latency, duty-cycle use, share of non-severe events escalated, duty violations (must be 0). Policy parameters are the stored tuned values; nothing is tuned on the held-out grid.
\textbf{G2 layer attribution} (one mechanism at a time against the baseline P1+, then a 2x2 or 2x2x2 factorial for interactions, both grids; intervals from seeds): PHY (SF per class: Critical one SF higher; coding rate 4/5 to 4/8), MAC (retry schedule: immediate vs spaced 5 s; hold on/off; backoff; reserved duty share 0 vs 0.35; capture on/off), network and security (tag length 4, 8, 16; signed ACK on/off meaning ACKs of 4 bytes less and spoofable [airtime effect only]; CANCEL on/off), application (class rule point vs upper end; interval upper end level; abstain routing). Interaction: SF raise x retries, hold x spacing, tag length x SF, each as a full factorial. Reported: which layer carries the gain and which layers do not matter. DESIGN-grid results are not reported as confirmed.
\textbf{G3 mechanisms}, each against the best earlier policy on the HELD-OUT grid; the baseline policy is the policy with the highest mean critical reach within 10 s on the DESIGN grid (decided before the held-out run, recorded in the table):
\begin{itemize}
\item \textbf{H-G3a (slotted fast path against pure ALOHA).} Critical reach within 10 s at least 5 points higher (interval excluding zero) in at least 60\% of the held-out cells (at least 10 of 16), AND no cell lost by more than 2 points. \textbf{Counts against:} fewer than 60\% of cells or any cell lost by more than 2 points. Beacon cost reported as airtime and as share of the gateway's duty (assumed).
\item \textbf{H-G3b (listen-before-talk against duty-cycle-only access).} Same criterion.
\item \textbf{H-G3c (V2V relay against no relay).} Critical reach within 30 s (reception at the gateway) at least 5 points higher (interval excluding zero) in at least 60\% of the held-out cells with p >= 0.30 (all 8 such cells in the grid), AND none of those cells lost by more than 2 points. Extra airtime and the duplicate-alert rate (extra gateway receptions per critical alert) reported.
\end{itemize}
\textbf{G4 regime map} over (density, loss, SF, duty) from G1 and G3: the policy with the best critical reach within 10 s per cell with margin and interval; ties where the best beats the runner-up by less than its interval. HELD-OUT cells give the map; DESIGN cells are shown, marked. Where P3 wins, if anywhere, is reported.
\textbf{G5 sensitivity} (HELD-OUT cells, 30 seeds): one assumption at a time, stated amounts: shadowing sigma 6 -> 3 and 9 dB; path-loss exponent 3.0 -> 2.5 and 3.5; capture threshold 6 -> 3 and 10 dB; ACK delay 0.1 -> 0.5 and 1.0 s; tag length 8 -> 4 and 16; relay delay U(1,5) -> U(0.5,2) and U(3,8) s; hear probability 0.5 -> 0.25 and 0.75; slot count 8 -> 4 and 16; sensing probability 0.7 -> 0.4 and 0.9; channels 1 -> 3; gateways 1 -> 2; severity mix (range +-16 -> +-32 g pairs). Reported: change in the main result (critical reach within 10 s, per policy and mechanism) and whether any conclusion flips; the conclusions that depend on a single assumption are listed.
\textbf{G6 regression:} the 48 attack cases (+ 9 designed cases) and the bounded model check of B1 are re-run against every protocol change used in G2 and G3 (tag length 4 and 16; relay (duplicate re-ACK budget); slotted fast path and LBT (timing only: no change to message rules; stated); CANCEL variants: design, pending cancel). New counterexamples or failing cases are reported. A check of the model and the designed cases, not the implementation.
\textbf{G7 summary:} one row per finding: layer, claim, evidence label SIMULATED, design-grid result, held-out result, sensitivity verdict, new / confirmation / negative. No DESIGN-grid result is a confirmed finding. Negative layer findings go into \texttt{docs/03\_results/NEGATIVES\_REGISTER.md}. The simulator is not tuned to produce a win; ties and losses are reported.

\subsection*{Summary table}

{\scriptsize\raggedright
\begin{longtable}{p{0.190\linewidth}p{0.190\linewidth}p{0.190\linewidth}p{0.190\linewidth}p{0.190\linewidth}}
\toprule
Study & Seeds & Baseline & Primary metric & Origin \\
\midrule
\endhead
B2 & 5000 + h (30) & P0 & critical reach within 10 s & SIMULATED \\
B3 & n/a & Semtech formula; 1 - exp(-2G) & airtime error; collision probability error & ANALYTICAL, SIMULATED \\
C3 & 42 + k & real-data smallest range & required range & SIMULATED \\
F1 & 0 to 4; 42 & original estimators and criteria & per hypothesis & MEASURED\_SOFTWARE, EXPLORATORY-REVISED \\
F2 & 0 to 4 & real-only training & recall at 1\% target; AP; decision time & MEASURED\_SOFTWARE, EXPLORATORY-REVISED \\
F3 & 9000 + h (30) & lowered-threshold table & mis-triage; critical reach within 10 s & SIMULATED \\
G & 6000 + h, 8000 + h (30) & best earlier policy & critical reach within 10 s (30 s for relay) & SIMULATED \\
\bottomrule
\end{longtable}}

\subsection*{Addenda}

\subsubsection*{Addendum 1 (2026-11-13; after B1, B2 and B3 were run; nothing in Blocks C, F or G had been run)}

\begin{itemize}
\item \textbf{B1.} The watermark variant was not modelled, so the sentence in B1 that a watermark "would cancel a delayed earlier event" is a statement from reasoning (inferred), not a model-check result. The exploratory two-event bound that was run (EXT2: counters up to 4, 2 events, 2 alert sends, 1 flash, 2 attacker actions, 1 lost message) covers the exact-counter pending-cancel change and the design only. The 48 original attack cases use \texttt{BaseVerifier} (no CANCEL) and are unaffected by construction; the designed cases of \texttt{attacks\_extra} use \texttt{SctBase}.
\item \textbf{B3.} The Semtech document could not be opened; (i) is a test against the formula as transcribed and hand arithmetic (stated in DECISIONS). The sensitivities of A-G4 stay "assumed".
\item \textbf{B2.} No deviation. The class pools of P1+e are the lowered pools of \texttt{t1} (f of P3's escalation), as written.
\end{itemize}

\subsubsection*{Addendum 2 (2026-11-13; written after Blocks B, C and G were run, and after the sealed F1 run; before the F1 cross-validation runs, the F2 test scoring and the F3 runs were evaluated)}

\begin{itemize}
\item \textbf{C generator.} The generator design changed while its fidelity was inspected on the same real records (in-sample), before any C3, F or G use: Gaussian copula with empirical marginals instead of a multivariate normal; a plateau-plus-spike blend to reach the real area ratios; calibration of the peak and duration medians only (delta-V not shifted). These are generator design choices, not hypothesis changes. Result: only frontal barrier and other pass; pole and side fail and are excluded from C3, F and G (the preregistered gate). H-C3 is therefore tested on frontal barrier only (the pole mode is NOT tested).
\item \textbf{F1.} H-F1d (original H-X6 on splits G and O) and H-F1e (original H-X5) are NOT\_RUN: the 176-test grouped-split machinery was not re-implemented with synthetic training inside the folds in the available time. Their original verdicts stand. Candidates for H-F1 are the augmented direct GBM and the augmented 1D CNN, mode none, as written; the anchored estimator and the reliability score use real rows plus synthetic rows without the augmented copies (as the original models did). Only the +-16 g cell enters H-F1b and H-F1c.
\item \textbf{F2.} The main estimate on the VZCrash side is the supplied-split test result (de-duplicated primary, strict sensitivity); grouped nested cross-validation is NOT\_RUN there (no vehicle identifier). H-F2c is NOT\_RUN as preregistered. H-F2e (six length-specific forests of N19) and H-F2f (leave-one-category-out calibration variants V0 to V5) are NOT\_RUN in the available time. H-F2d reports the AP clause only; the flash clause is arithmetic.
\item \textbf{F0.} The surrogate generator has no numeric pass threshold in the preregistration. It is used as it is, with the marginals that match (KS at most 0.20) and those that do not reported in \texttt{f.f2.surrogate.ks.*}; the feature \texttt{post\_quiet} does not match.
\item \textbf{F3.} The pool-level clause uses the out-of-fold pairs of the F1 cross-validation (main estimate); the simulated clause compares P3 with the equal-escalation SCT policy P2e (and P1e), as the original N28 rule did.
\item \textbf{G.} The DESIGN grid was used to choose only the baseline policy of G3 (P4). The relay hypothesis uses reception at the gateway within 30 s (\texttt{reach30\_rx\_critical}) in the 8 cells with loss 40\%. The slot beacon is charged to the gateway's duty allowance; the share is reported against the allowance of the 5,400 s that are simulated.
\item \textbf{Holm.} NOT\_RUN hypotheses enter the family with p = 1 (family size 18 unchanged).
\end{itemize}

\subsubsection*{Addendum 3 (2026-11-13; correction, written when the full test suite flagged the count)}
The text of B1 says "48 + 9 attack cases"; the designed extra cases of \texttt{attacks\_extra} number 10 (X01 to X10), not 9. Every result of this round uses the real count (10); the sentence above is not edited. \texttt{tests/test\_consistency.py} treats the preregistration like \texttt{docs/07\_log/PROGRESS.md} (a dated record).

